\section{Problem Definition and Motivation}
\label{sec:problem}

% 证据：exact/mirror.py（模型、目标）；exact/family_bounds.py、robot_lb.py（下界）；tools/gen_family*.py、gen_request.py、rulegen/instances.py（实例与划分）；
%       实验0_表示分析/exp0lib.py（D0～D5）；实验0_报告.md、out/report_metrics.md、out/report_writeback.md（III-E 的数字）；排班_解码方案汇总.md（最优解结构）
% 2026-10-22 已两轮独立核对；本版只改写表达，事实未改。

This section defines the rostering problem, the instances, and the six decoders that we use as baselines.
It ends with an analysis of these decoders, which shows what a better decoder needs.

\subsection{Robot--Tool Duty Rostering}
\label{sec:problem:model}

A duty post must be staffed by one robot at all times during a horizon $[0,\horizon]$.
Robots take turns in \emph{shifts}, and each shift starts when the previous one ends.
If shift $i-1$ ends at time $E_{i-1}$ and shift $i$ is served by a robot of type $\tau$,
shift $i$ ends at $E_i = E_{i-1} + w_\tau$, where $w_\tau$ is the duty length of type $\tau$.
The first shift starts once the first robot has installed its tools and travelled to the post.
The roster is complete when $E_i \ge \horizon$.

\textbf{Robots.}
Each robot type $\tau$ has a travel time to the post $o_\tau$, a return time $b_\tau$, a charging time $\gamma_\tau$ and a maintenance time $\mu_\tau$.
A robot can serve at most $n_\tau$ duties between two maintenances; we call this its \emph{duty budget}.
After each duty the robot returns and its tools are removed.
It then charges, or, if its budget is used up, goes to maintenance, which refills the budget.
A robot is \emph{available} for a shift that starts at time $E$ if it is ready by $E - o_\tau$ and still has budget for one duty.
At most $L_\tau$ robots of type $\tau$ may be used.

\textbf{Tools.}
Each duty needs $\nu_\tau$ tools of types compatible with the robot.
A tool in use is busy for its installation, its travel and work time, its removal and a rest period.
When a robot is dispatched, the simulator first reuses compatible tools that are ready, earliest first.
If tools are missing, new ones are created, and their type has to be chosen.
At most $L_k$ tools of type $k$ may be used.
A roster is infeasible if it needs more robots or tools than these limits allow.

\textbf{Decisions.}
At every shift $i$ one decision is made: which robot serves it.
The decision is an action
\begin{equation}
  \act_i \in \{\reuse(r) : r \in \avail_i\} \cup \{\new(\tau)\},
\end{equation}
where $\avail_i$ is the set of available robots.
\reuse{}$(r)$ sends a robot that is already in the fleet; \new{}$(\tau)$ adds a new robot of type $\tau$.
Together with the types of the new tools, these actions fix the roster.
Each action changes which robots are available later, so a decision at one shift affects all later shifts.

\subsection{Objectives}
\label{sec:problem:objectives}

We minimize two objectives.
The first is the \emph{resource count}, the number of robots and tools used:
\begin{equation}
  \fres = |\mathcal{R}| + |\mathcal{T}|.
\end{equation}
The second is the \emph{imbalance} of the robot fleet:
\begin{equation}
  \fbal = 0.3\,\sigma(\mathbf{n}) + 0.4\,\sigma(\mathbf{h}) + 0.3\,\sigma(\Delta\mathbf{t}),
\end{equation}
where $\sigma(\cdot)$ is the standard deviation, $\mathbf{n}$ and $\mathbf{h}$ are the robots' duty counts and total duty hours,
and $\Delta\mathbf{t}$ are the gaps between the times at which successive robots first go to maintenance.
A low imbalance means that all robots wear evenly and do not go to maintenance at the same time.
The two objectives conflict.
A small fleet must work some robots hard, while an even workload needs extra robots.

For each instance we also compute a lower bound $\LB$ on the resource count.
The robot part comes from an exact feasibility search with unlimited tools;
the tool part checks whether a tool used in one shift can be ready again for the next shift, and if not, counts the extra tools needed.

\subsection{Instances}
\label{sec:problem:instances}

We use 62 instances (Table~\ref{tab:instances}).
Four are hand-built scenarios: one is the request file of the scheduling software that motivated this work (one robot type, 700\,h),
and three have three robot types and four tool types, one of them with tight fleet limits.
The other 58 are random instances in three families:
\begin{itemize}
  \item \textbf{E} (homogeneous): one robot type.
  \item \textbf{H} (heterogeneous): two or three robot types, all with 10\,h duties.
  \item \textbf{M} (mixed duty lengths): two or three robot types with duty lengths from $\{6,8,10,12\}$\,h, so shifts have different lengths.
\end{itemize}
Robot parameters are drawn at random: travel 1--2\,h, charging 2--5\,h, maintenance 48--120\,h,
2--5 duties per maintenance cycle, 2 or 3 tools per duty, and 6--10 robots per type.
Tools take 1--2\,h to install or remove and rest for 3--8\,h.
We generated 60 random instances and removed two for which no method found a feasible roster.

The instances are split once into 15 training instances (five per family) and 47 test instances.
The construction rules of our method are learned on the training instances only, and its deviation density is selected on them (Section~\ref{sec:method}).
The test instances measure the final performance.
% 待定（κ 完整结果出来后再写）：是否说明早期探索性运行曾在测试集上评价过 κ = 0.5

\begin{table}[t]
  \centering
  \caption{Instance groups. $\Nmax$ is the upper bound on the number of shifts.}
  \label{tab:instances}
  \footnotesize
  \setlength{\tabcolsep}{2.5pt}
  \begin{tabular}{@{}lccccccc@{}}
    \toprule
    Group & Train & Test & \shortstack{Robot\\types} & \shortstack{Tool\\types} & Horizon (h) & Duty (h) & $\Nmax$ \\
    \midrule
    Hand-built & 0 & 4  & 1--3 & 1--4 & 700--1440 & 6, 10 & 70--240 \\
    E        & 5 & 12 & 1    & 1--3 & 480--1080 & 10    & 48--108 \\
    H        & 5 & 15 & 2--3 & 2--3 & 480--1080 & 10    & 48--108 \\
    M        & 5 & 16 & 2--3 & 2--3 & 480--1080 & 6--12 & 60--180 \\
    \midrule
    Total    & 15 & 47 & & & & & \\
    \bottomrule
  \end{tabular}
\end{table}

\subsection{Encoding and Baseline Decoders}
\label{sec:problem:decoders}

All decoders in this paper use a random-key genotype $\genome \in [0,1)^n$, a vector of $n$ numbers, and build the roster one shift at a time.
At each shift the decoder finds the available robots, picks an action, and simulates the shift.
Since every shift lasts at least $\min_\tau w_\tau$, the number of shifts is at most
\begin{equation}
  \Nmax = \left\lceil \horizon / \min_\tau w_\tau \right\rceil .
\end{equation}

The baseline decoders give each shift two genes: a robot gene $x_i$ that picks the action, and a tool gene $y_i$ that picks the type of any new tool.
The tool gene points into the list of the $c_\tau$ compatible types at position $\lfloor y_i (c_\tau - 1) \rfloor$;
if that type has reached its limit, the next type is used.
All decoders read the tool genes in this way, following the original software.
They differ in how they read the robot gene.
We name them \dec{0} to \dec{5}; each one adds a piece of problem knowledge to the one before.

\textbf{\dec{0} (original).}
The gene is scaled to an index over the available robots followed by the robot types: $\lfloor (x_i - 10^{-3})(|\avail_i| + n_R) \rfloor$, where $n_R$ is the number of types.
A small index sends an available robot, and a large one opens a new robot.
Because $|\avail_i|$ changes from shift to shift, the same gene value can mean ``send'' at one shift and ``open'' at another.

\textbf{\dec{1} (threshold).}
A gene below $\theta = 0.85$ always means \reuse, and a gene above $\theta$ always means \new.
Below $\theta$, the gene picks robot $\lfloor x_i |\avail_i| / \theta \rfloor$; above $\theta$, it picks a type.

\textbf{\dec{2} (ranking).}
As \dec{1}, but the available robots are sorted, fewest duties first (ties: most remaining budget first).
A small gene now sends the least-used robot.

\textbf{\dec{3} (capacity look-ahead).}
As \dec{2}, but before opening a new robot the decoder estimates how many duty hours the current fleet can still provide.
If the fleet can cover the rest of the horizon, it sends an existing robot instead.

\textbf{\dec{4} (maintenance phase).}
As \dec{3}, but robots that are closest to maintenance are ranked first.

\textbf{\dec{5} (backtracking repair).}
As \dec{3}, with one more step.
If no robot is available at shift $i$, the decoder goes back over the last $W = 12$ shifts.
At each of them it tries sending a different robot and re-simulates up to shift $i$.
It keeps the first change that leaves a robot available at shift $i$, and opens a new robot only if no change works.

\dec{5} performs best of the six and is our main baseline.
All six decoders, however, share one structure: every shift has its own gene, and that gene decides the shift.
We implemented all decoders in Python and checked \dec{1}--\dec{5} against the C++ code of the original software, genotype by genotype.

\subsection{Motivation: What the Baseline Decoders Show}
\label{sec:motivation}

To see why some decoders work better than others, we measured four properties of each decoder on 16 instances
(the four hand-built scenarios with all six decoders, and 12 E, H and M instances with \dec{2}--\dec{5}):
\begin{itemize}
  \item \textbf{Neutral rate}: the chance that redrawing one robot gene leaves the roster unchanged.
  A high neutral rate means that many genotypes produce the same roster, \ie the encoding is \emph{redundant}.
  \item \textbf{Good-start rate}: the chance that a uniformly random genotype gives a roster with $\fres \le \LB + 1$.
  \item \textbf{Change propagation}: after one robot gene is changed, the share of later shifts whose robot also changes.
  \item \textbf{Reachability}: whether some genotype produces the optimal roster. We computed the optimal rosters of the hand-built scenarios by dynamic programming.
\end{itemize}
We also recorded the \emph{override share}: how often the capacity look-ahead of \dec{3}--\dec{5} (below) turns a \new action into a \reuse action.
The good-start rate is estimated from 1000 random genotypes.
The neutral rate and change propagation are measured on 100 random feasible genotypes and on the 30 best of the 1000, with five single-gene changes each.
We then compared these properties with the hypervolume (HV) that NSGA-II \cite{deb2002nsga2} reaches with each decoder
(population 30, about 465 evaluations).
For each instance we computed the Spearman correlation between each property and HV across the decoders, and we report the median over instances.
Table~\ref{tab:neutral} lists the first two properties on the hand-built scenarios.

\begin{table}[t]
  \centering
  \caption{Neutral rate / good-start rate of the six decoders on the four hand-built scenarios.}
  \label{tab:neutral}
  \footnotesize
  \setlength{\tabcolsep}{1.6pt}
  \begin{tabular}{@{}lcccccc@{}}
    \toprule
    Scenario & \dec{0} & \dec{1} & \dec{2} & \dec{3} & \dec{4} & \dec{5} \\
    \midrule
    1 type, 700 h   & 0.30/.000 & 0.42/.007 & 0.41/.007 & 0.61/.919 & 0.61/.929 & 0.69/.981 \\
    3 types, 720 h  & 0.10/.000 & 0.21/.000 & 0.22/.000 & 0.44/.020 & 0.46/.024 & 0.50/.032 \\
    3 types, 1440 h & 0.06/.000 & 0.13/.000 & 0.14/.000 & 0.38/.014 & 0.42/.008 & 0.44/.029 \\
    3 types, tight  & 0.12/.000 & 0.24/.000 & 0.19/.000 & 0.44/.020 & 0.47/.024 & 0.49/.032 \\
    \bottomrule
  \end{tabular}
\end{table}

\textbf{Finding 1: the better decoders favor good rosters.}
From \dec{0} to \dec{5}, the neutral rate rises from 0.06--0.30 to 0.44--0.69, and the good-start rate rises with it.
Both go with better HV: the median correlation is $+0.91$ for the neutral rate and $+0.80$ for the good-start rate,
with the same sign on every instance where the correlation is defined (12/12 and 14/14).
The override share correlates even more strongly ($+0.97$, 16/16).
These three properties rise together, so the data cannot tell their effects apart.
Taken together, they agree with the theory of redundant representations \cite{rothlauf2006representations}:
an encoding helps when its redundancy favors good solutions.

\textbf{Finding 2: change propagation is high in every decoder.}
On the hand-built scenarios, changing one gene changes the robot at 39--61\% of the later shifts, for every decoder.
This share does not track performance (median correlation $+0.23$, same sign on 7/12 instances).
The reason is structural: a changed decision changes which robots are available, and every later gene is then read in a different situation.
None of the steps from \dec{0} to \dec{5} reduced this effect.

\textbf{Finding 3: built-in knowledge helps even when it hides good rosters.}
The capacity look-ahead of \dec{3}--\dec{5} sometimes forbids opening a robot.
On two hand-built scenarios, this makes the optimal roster unreachable: that roster opens a fifth robot at shift 9, which the look-ahead does not allow.
Other optimal rosters can still be reached, and on one of the two scenarios \dec{5} still finds the optimal values.
Making the optimal roster reachable again did not help.
When we added a range of gene values that always opens a robot, the imbalance got significantly worse on the two scenarios that had not lost the optimum,
and nothing improved on the two that had.

\textbf{Finding 4: well-balanced rosters are rare.}
With every decoder, at most 1.7\% of random genotypes give a roster with $\fres = \LB$ and an imbalance within 1.5 times the optimum;
on the three-type scenarios almost none do.
This is where the search falls short: at $\fres = \LB$, the mean imbalance reached by \dec{5} is 1.1--2.3 times the optimum.

\textbf{Structure of optimal rosters.}
The optimal rosters follow a simple pattern.
In 85--90\% of their shifts they send the available robot with the fewest duties.
In the remaining shifts they send a different robot; in the single-type scenario, this happens in 10 of its 65 reuse decisions.
A per-shift encoding has no direct way to say ``follow this rule except at a few shifts''.
The choice of rule also matters.
On the four hand-built scenarios, always sending the robot with the fewest duties needs 11--17 resources,
while always sending the robot with the least remaining budget reaches the lower bound of 9 on all four.

\textbf{What a better decoder needs.}
These results lead to three requirements:
\begin{enumerate}
  \item The encoding should express ``a rule plus a few exceptions'' directly: the genes should say only where and how to depart from the rule.
  \item How strongly the decoder favors the rule should be adjustable, because the distance between the rule and good rosters differs between rules and instances.
  \item The rule should be chosen for the rosters that the search finds around it, not only for the single roster it builds.
\end{enumerate}
Section~\ref{sec:method} presents a decoder that meets these requirements.
